% ==================================================
% Motivation
% Discrete Calculus Notes Stub
% ==================================================

\section*{Motivation}

% TODO: Add definitions, theorems, and exposition here.

% ==================================================
% Motivation
% Discrete Calculus
% ==================================================

\section*{Motivation}

Classical calculus studies functions defined on continuous domains, typically
functions \(f : \mathbb{R} \to \mathbb{R}\).
Its central tools are the derivative and the integral:

\begin{itemize}
\item The derivative measures \emph{instantaneous change}.
\item The integral accumulates \emph{total change}.
\end{itemize}

However, many important mathematical objects are inherently \emph{discrete}.
Examples include

\begin{itemize}
\item sequences \(a_n\),
\item recurrence relations,
\item sums of integers or combinatorial quantities,
\item discrete probability distributions,
\item algorithms and computational processes.
\end{itemize}

In these contexts the domain is typically the natural numbers

\[
f : \mathbb{Z} \to \mathbb{R}
\]

rather than the real line.

To study change in such settings, calculus is replaced by its discrete
counterpart.

\subsection*{Discrete Analogs of Calculus}

The central idea of \emph{discrete calculus} is that the roles of derivative and
integral are played by two simple operations.

\begin{center}
\begin{tabular}{c|c}
Continuous Calculus & Discrete Calculus \\
\hline
Derivative \( \frac{d}{dx} f(x) \) & Forward difference \( \Delta f(n) \) \\
Integral \( \int f(x)\,dx \) & Summation \( \sum f(n) \)
\end{tabular}
\end{center}

The forward difference operator measures change between consecutive values of a
function.

\[
\Delta f(n) = f(n+1) - f(n).
\]

This quantity plays the same role for sequences that the derivative plays for
functions of a real variable.

Likewise, summation serves as the discrete analogue of integration.

\[
\sum_{k=a}^{b} f(k)
\]

accumulates the values of a function across a finite discrete interval.

\subsection*{Telescoping Phenomena}

One of the most important identities in discrete mathematics arises when
differences are summed.

\[
\sum_{k=a}^{b-1} \bigl(f(k+1)-f(k)\bigr)
=
f(b)-f(a).
\]

The intermediate terms cancel, leaving only the endpoints. Such sums are called
\emph{telescoping sums}.

This identity is the discrete counterpart of the
\emph{Fundamental Theorem of Calculus}.

\subsection*{Why Discrete Calculus Matters}

Discrete calculus appears throughout mathematics:

\begin{itemize}
\item in combinatorics and enumerative formulas,
\item in recurrence relations and generating functions,
\item in probability theory and stochastic processes,
\item in numerical analysis and finite difference methods,
\item in computer science and algorithm analysis.
\end{itemize}

For students of analysis, discrete calculus is particularly valuable because it
reveals the structural skeleton behind many identities involving sums and
sequences.

Many techniques used to estimate limits of sequences or series ultimately rely
on discrete analogues of differential identities.

\subsection*{A Structural View}

Continuous calculus studies how functions change along the real line.

Discrete calculus studies how functions change along the integers.

Despite the difference in setting, the two theories mirror each other
remarkably closely. Understanding this parallel provides powerful intuition and
useful techniques for manipulating sums, sequences, and recurrence relations.

In the sections that follow we develop the basic tools of discrete calculus,
beginning with the forward difference operator and the algebra of finite
differences.

% ==================================================
% Discrete vs Continuous Calculus Correspondence
% ==================================================
\begin{center}
\textbf{Structural Correspondence Between Continuous and Discrete Calculus}
\end{center}

\begin{center}
\begin{tikzpicture}[
    node distance=2.8cm,
    every node/.style={rectangle,draw,rounded corners,align=center,minimum width=4.2cm,minimum height=1cm},
    arrow/.style={->,thick}
]

% Continuous column
\node (deriv) {$\displaystyle \frac{d}{dx} f(x)$\\Derivative};
\node (int) [below of=deriv] {$\displaystyle \int f(x)\,dx$\\Integral};
\node (power) [below of=int] {$\displaystyle \frac{d}{dx} x^n = n x^{n-1}$\\Power Rule};
\node (ftc) [below of=power] {Fundamental Theorem\\of Calculus};

% Discrete column
\node (diff) [right=5cm of deriv] {$\Delta f(n)=f(n+1)-f(n)$\\Forward Difference};
\node (sum) [below of=diff] {$\displaystyle \sum f(n)$\\Summation};
\node (falling) [below of=sum] {$\Delta n^{\underline{k}} = k\, n^{\underline{k-1}}$\\Falling Factorial Rule};
\node (telescoping) [below of=falling] {Telescoping Sum\\Identity};

% arrows
\draw[arrow] (deriv) -- (diff);
\draw[arrow] (int) -- (sum);
\draw[arrow] (power) -- (falling);
\draw[arrow] (ftc) -- (telescoping);

\end{tikzpicture}
\end{center}



% ============================================================
% CONTINUOUS VS FINITE CALCULUS
% Structural dictionary between differential calculus
% and discrete (finite) calculus
% ============================================================

\section*{Continuous vs Finite Calculus}

\begin{center}
\renewcommand{\arraystretch}{1.6}
\begin{tabular}{|c|c|}
\hline
\textbf{Continuous Calculus} & \textbf{Finite (Discrete) Calculus} \\
\hline
Function $f(x)$ & Function $f(n)$ \\
\hline
Derivative operator $\dfrac{d}{dx}$ & Forward difference $\Delta f(n)=f(n+1)-f(n)$ \\
\hline
Higher derivatives $\dfrac{d^k}{dx^k}$ & Higher differences $\Delta^k f(n)$ \\
\hline
Integral $\displaystyle \int_a^b f(x)\,dx$ & Sum $\displaystyle \sum_{k=a}^{b-1} f(k)$ \\
\hline
Antiderivative $F'(x)=f(x)$ & Discrete antiderivative $\Delta F(n)=f(n)$ \\
\hline
Fundamental theorem of calculus &
Fundamental theorem of finite calculus \\
&
$\displaystyle \sum_{k=a}^{b-1} \Delta f(k)=f(b)-f(a)$ \\
\hline
Power function $x^n$ & Falling power $x^{\underline{n}}$ \\
\hline
Power rule &
Finite power rule \\
$\dfrac{d}{dx}x^n = n x^{n-1}$ &
$\Delta x^{\underline{n}} = n x^{\underline{n-1}}$ \\
\hline
Taylor series &
Newton series \\
\hline
Exponential $e^{ax}$ &
Discrete exponential $(1+a)^n$ \\
\hline
Product rule &
Discrete product rule \\
$\dfrac{d}{dx}(fg)=f'g+fg'$ &
$\Delta(fg)=f(n+1)\Delta g(n)+g(n)\Delta f(n)$ \\
\hline
Integration by parts &
Summation by parts \\
\hline
$\displaystyle \int f\,dg$ &
$\displaystyle \sum f\,\Delta g$ \\
\hline
Riemann sums &
Finite sums (exact) \\
\hline
Differential equation &
Difference equation \\
\hline
\end{tabular}
\end{center}

\vspace{1em}

\noindent
\textbf{Conceptual Parallel}

\[
\begin{aligned}
\text{Derivative} &\leftrightarrow \text{Difference} \\
\text{Integral} &\leftrightarrow \text{Summation} \\
\text{Power } x^n &\leftrightarrow \text{Falling power } x^{\underline{n}} \\
\text{Taylor series} &\leftrightarrow \text{Newton series}
\end{aligned}
\]

\vspace{1em}

\noindent
\textbf{The Central Identity (Discrete FTC)}

\[
\sum_{k=a}^{b-1} (f(k+1)-f(k)) = f(b)-f(a)
\]

which mirrors

\[
\int_a^b f'(x)\,dx = f(b)-f(a).
\]
